%!TeX root = main.tex
%!encoding = UTF-8 Unicode
\section{State of the Art and Research Gap in Semiconductor-Focused Emission Characterization}

As established in Chapter~\ref{ch:Fundamentals}, semiconductor switching generates the voltage and current transients that excite EM emissions, while the surrounding circuit determines their propagation to system-level measurement ports.
This distinction is also reflected in the literature on semiconductor-focused EMC analysis~\cite{hillenbrand_frequency_2016,takahashi_noise-source_2021,marlier_modeling_2012}.
Characterizing a device therefore requires distinguishing its switching-source spectra from the effects of the propagation path.
Here, intrinsic emission behavior denotes the terminal-level switching behavior under defined operating and circuit conditions, rather than a device property independent of its environment.

The present review focuses on approaches that relate EM emissions to semiconductor switching behavior, particularly the high-frequency switching behavior not captured by idealized pulse models~\cite{kragl_impact_2024}.
Three approaches are considered: system-level measurements, simulation-based analysis, and direct measurements at the device terminals.
They are assessed in terms of their ability to distinguish individual source contributions, resolve switching events, and characterize small high-frequency signal components.

\subsubsection*{System-Level Measurements}
In experimental system-level setups, various semiconductor technologies are tested within a commutation cell, and their emissions are compared using a LISN~\cite{apelsmeier_model_2020, takahashi_noise-source_2021, zhang_two-port_2020, tlig_conducted_2025, lemmon_methodology_2018, huber_ultra-low_2017} or antenna~\cite{radchenko_transfer_2014,domurat-linde_analysis_2012}. 
This approach enables comparisons of the system-level emissions obtained with different semiconductor technologies and provides a direct link to conventional EMC assessment.

However, the measured spectrum contains the combined contributions of all active noise sources after propagation through the surrounding circuit.
It does not, by itself, identify the contribution of an individual switch or switching event, which makes a detailed analysis challenging.
Additionally, within a given frequency range, the combined emissions can be dominated by a single switching event whose contribution at the measurement port substantially exceeds those of the other events.
If a parameter variation reduces the emissions associated with a weaker event while leaving the dominant contribution unchanged, the resulting change in the total level may be too small to distinguish in the logarithmic emission spectrum.
A beneficial effect on that event may therefore remain undetected in a system-level measurement, despite a substantial reduction in its individual contribution.

Furthermore, the measured emissions depend on the passive components and layout through their transfer functions~\cite{xie_high_2019,phukan_characterization_2023}.
Strong attenuation for a given frequency range along the propagation path can mask parameter-induced changes in the source spectrum.
A weak dependence of the measured emissions on a device parameter may therefore be misinterpreted as low sensitivity of the source itself to that parameter.
To account for the propagation-path influence, studies have reconstructed noise-source spectra from system-level measurements by deconvolving the propagation-path response.
While this approach provides satisfactory results at low frequencies, its accuracy in the higher-frequency range governed by the detailed switching transients remains insufficiently established.

\subsubsection*{Simulation-Based Approaches}
The limitations of system-level measurements make it difficult to determine which aspects of semiconductor switching behavior are responsible for a change in the measured emissions.
Simulation-based approaches provide both direct access to true device voltages and currents, and can model the propagation path to also predict the resulting system-level emissions, enabling a more comprehensive understanding of the relationship between device behavior and observed emissions.
In these studies, behavioral models are parameterized to identify which device properties drive emissions in specific frequency bands~\cite{hillenbrand_sensitivity_2017, kragl_semiconductor-focused_2025, liu_modeling_2016, yuan_emi_2015}.
Simulation permits individual device parameters and source contributions to be investigated separately, subject to the validity of the model and its representation of circuit parasitics.

However, these models are calibrated using conventional double-pulse tests and static measurements~\cite{van_der_broeck_simulating_2023,kragl_semiconductor-focused_2025}.
The resulting calibration data do not capture the small high-frequency signal components relevant to EMC when these fall below the measurement noise floor, as explained in Section~\ref{sec:MeasurementLimitation}.
Without this information, the models cannot be assumed to predict these components accurately.
Their behavior in this frequency range remains uncalibrated and unverified.
Accurate high-frequency voltage and current measurements are therefore needed to characterize the devices and validate the models.

\subsubsection*{Direct Measurements at Device Terminals}

Direct measurements at the device terminals provide the voltage and current signals needed for device characterization and model validation.
They allow the source spectra to be investigated without reconstructing them from system-level emissions or relying on a semiconductor model.

The main difficulty is resolving small high-frequency components within a large switching transient.
Conventional voltage measurements can become noise-limited above the device resonance frequency despite sufficient probe bandwidth~\cite{labrousse_switching_2011,kragl_impact_2024}.
Sufficient signal-to-noise ratio (SNR), rather than nominal bandwidth alone, is therefore essential.

To address this limitation,~\cite{oswald_high-bandwidth_2011} combined a passive voltage probe with a high-pass filter to measure high-frequency switching components in a discrete IGBT\@.
However, the approach was demonstrated only for IGBTs and for frequencies up to \SI{30}{MHz}.
The passive probe measured only the low-side drain-source voltage, providing access to just one of the three source signals in a switching cell: the high-side voltage, low-side voltage, and commutation current.

\subsubsection*{Research Gap and Methodological Requirements}
The remaining gap is a validated method for measuring all three source signals of a switching cell with sufficient SNR up to \SI{400}{MHz}.
These signals are the high-side voltage, low-side voltage, and commutation current.
The method must also separate the spectral contributions of individual switching events.
Measurements under defined operating and circuit conditions can then support device comparisons and model validation.

This thesis addresses this gap through measurement methods, sensors, and signal processing.
Section~\ref{sec:MeasurementLimitation} first explains the noise limits of conventional measurements.
The rest of this chapter develops methods to reduce the effects of measurement noise and separate switching events.
Chapter~\ref{ch:Sensors} presents the sensors, and Chapter~\ref{ch:Validation} validates the complete measurement setup.



